\documentclass[12pt]{article}
\usepackage{amsmath}
\usepackage{amssymb}
\usepackage{enumerate}

\setlength{\parskip}{0.2cm}
\setlength\parindent{0pt}

\begin{document}

\textbf{1. LATEX Warm-up}

Name: \textbf{Sanjida Akter}
Date: \textbf{02/18/2026} 
Assignment \#: \textbf{1}

\hrulefill

\textbf{2. Basic \LaTeX\ Practice}

The people of Cappa found that there are exactly
\[
\sqrt{\frac{n^5 - 3n}{\max\{\ln n, 0.4 \cdot \cos n\}}}
\]
trees on their planet, CappaHoma\textsubscript{25}, where $n$ is the number of Cappians.

\textbf{3. \LaTeX\ Practice: Sets}

The objects $S_1 = \{\text{Foo}, \text{Bar}, \text{Baz}\}$ and 
$S_2 = \langle \text{Foo}, \text{Bar}, \text{Baz} \rangle$ may seem to be the same, but they are not.

The first is a set, while the second is a sequence.

Agreeably, the set $S_1$ is not as famous as $\mathbb{N}$ or $\mathbb{R}$, but knowing that $Foo \in S_1$ is itself a great honor for $S_1$.

\textbf{4. Equality of Sets and Sequences}

\begin{enumerate}[i.]

\item 
Yes, the two sets $\{340,1.5,-32\}$ and $\{1.5,340,1.5,-32,-32\}$ are equal because sets ignore the order of elements and repeated elements, so both contain exactly $340$, $1.5$, and $-32$.

\item 
No, the two sets $\{340,1.5,-32\}$ and $\{1.5,340,32\}$ are not equal because $-32$ is not equal to $32$.

\item 
No, the sequences $\langle340,1.5,-32\rangle$ and $\langle1.5,340,1.5,-32,-32\rangle$ are not equal because sequences depend on the order and the number of elements. The first sequence has $3$ elements, while the second has $5$ elements.

\item 
No, the sequences $\langle340,1.5,-32\rangle$ and $\langle1.5,340,32\rangle$ are not equal because their first elements differ: $340$ is not equal to $1.5$.

\end{enumerate}

\textbf{5. Ways to Write Out Sets}

\begin{enumerate}[i.]

\item 
No, we cannot list all elements of $\mathbb{N}$ using the roster method because it contains infinitely many elements.

\item 
Three elements of 
$S = \{2n - 7 \mid n \in P, n \ge 15, n \text{ is odd}\}$ 
are $23$, $27$, and $31$, obtained by substituting $n = 15$, $17$, and $19$ into the expression $2n - 7$.

\item 
One possible set-rule notation for the cubes of even integers greater than or equal to $-10$ is
\[
\{n^3 \mid n \text{ is even},\ n \ge -10\}.
\]

\end{enumerate}

\textbf{6. Element Inclusion vs. Subsets}

Let $A = \{1,2,3,4,5\}$ and $B = \{2,4,6\}$.

\begin{enumerate}[i.]

\item 
(a) $3 \in A$ is true.  

(b) $6 \in B$ is true.  

(c) $1 \notin B$ is true.

\item 
The subsets of $B$ with exactly two elements are $\{2,4\}$, $\{2,6\}$, and $\{4,6\}$.

\item 
$B \subset A$ is false because $6 \in B$ but $6 \notin A$.

\item 
One possible set is $C = \{2,4,6,8\}$, which satisfies $B \subset C$ and $C \ne B$.

\end{enumerate}

\end{document}
